% Guía de usuario TeleCon4ESP32 — aplicación Android para makers.
% Compilar: pdflatex -output-directory=../build telecon_user_guide.tex
\documentclass[12pt,a4paper]{article}

\usepackage[margin=2.4cm]{geometry}
\usepackage[T1]{fontenc}
\usepackage[utf8]{inputenc}
\usepackage{lmodern}
\usepackage[spanish]{babel}
\usepackage{booktabs}
\usepackage{hyperref}
\usepackage{xcolor}
\usepackage{colortbl}
\usepackage{array}
\usepackage{enumitem}
\usepackage{needspace}
\usepackage{etoolbox}
\usepackage{tikz}
\usepackage{graphicx}
\usepackage{float}
\usepackage{caption}
\captionsetup{justification=centering,singlelinecheck=false}

\usetikzlibrary{arrows.meta,shadows}
\preto\section{\Needspace{9\baselineskip}}

\hypersetup{
  colorlinks=true,
  linkcolor=black,
  urlcolor=blue,
  pdftitle={Guía de usuario TeleCon4ESP32},
  pdfauthor={TeleCon}
}

\definecolor{ArduinoKeyword}{HTML}{00979C}
\definecolor{ArduinoFunc}{HTML}{D35400}
\definecolor{ArduinoComment}{HTML}{5D6D6E}
\definecolor{ArduinoPreproc}{HTML}{9B59B6}
\definecolor{TblHead}{HTML}{0F6F73}
\definecolor{TblSync}{HTML}{E8DAEF}
\definecolor{TblChan}{HTML}{D5F5E3}
\definecolor{TblKnob}{HTML}{D6EAF8}
\definecolor{TblSw}{HTML}{FCF3CF}
\definecolor{TblSum}{HTML}{FAD7A0}
\definecolor{TblEnd}{HTML}{E5E7E9}
\definecolor{TblTelem}{HTML}{D4E6F1}
\definecolor{TblId}{HTML}{F5CBA7}
\definecolor{WireApp}{HTML}{0F6F73}
\definecolor{WireBg}{HTML}{F4FBFC}

\newcommand{\thc}[1]{\textcolor{white}{\bfseries #1}}
\newcommand{\chip}[2]{\colorbox{#1}{\strut\,\scriptsize #2\,}}
\newcommand{\ui}[1]{\textsf{\textbf{#1}}}

\tikzset{
  telecon actor/.style={
    rounded corners=4pt, minimum width=2.7cm, minimum height=1.05cm,
    font=\sffamily\bfseries, text=white, inner sep=4pt
  },
  telecon step/.style={
    circle, inner sep=1pt, minimum size=0.52cm,
    font=\sffamily\bfseries\scriptsize, text=white
  },
  telecon pill/.style={
    rounded corners=2.5pt, inner sep=3.2pt,
    font=\sffamily\scriptsize, align=center
  },
  telecon arr/.style={-{Latex[length=2.1mm]}, line width=0.9pt}
}

% Prefer figures/es/<file>, then shared figures/<file>.
\newcommand{\shotlang}{es}
\newcommand{\shotimg}[1]{%
  \begin{tikzpicture}
    \node[inner sep=0] (img) {%
      \phantom{\includegraphics[width=0.68\textwidth,height=0.48\textheight,keepaspectratio]{#1}}%
    };
    \fill[black, fill opacity=0.12, rounded corners=13pt]
      ([xshift=-5.5pt+1.8mm,yshift=-5.5pt-2.3mm]img.south west) rectangle
      ([xshift=5.5pt+1.8mm,yshift=5.5pt-2.3mm]img.north east);
    \fill[black!92, rounded corners=13pt]
      ([xshift=-5.5pt,yshift=-5.5pt]img.south west) rectangle
      ([xshift=5.5pt,yshift=5.5pt]img.north east);
    \begin{scope}
      \clip[rounded corners=10pt]
        (img.south west) rectangle (img.north east);
      \node[inner sep=0] at (img.center) {%
        \includegraphics[width=0.68\textwidth,height=0.48\textheight,keepaspectratio]{#1}%
      };
    \end{scope}
    \draw[black!30, line width=0.65pt, rounded corners=13pt]
      ([xshift=-5.5pt,yshift=-5.5pt]img.south west) rectangle
      ([xshift=5.5pt,yshift=5.5pt]img.north east);
    \draw[TblHead, line width=1.25pt, line cap=round]
      ([xshift=-5.5pt+11pt,yshift=-5.5pt+3.4pt]img.south west) --
      ([xshift=5.5pt-11pt,yshift=-5.5pt+3.4pt]img.south east);
  \end{tikzpicture}%
}
\newcommand{\shot}[4]{%
  \par\addvspace{0.25em}%
  \IfFileExists{../../figures/\shotlang/#2}{%
    \begin{figure}[H]
      \centering
      \shotimg{../../figures/\shotlang/#2}%
      \caption{\textbf{#1.} #4\\[0.35em]
        {\sffamily\scriptsize\color{ArduinoComment}#3
          \enspace\textbullet\enspace\texttt{\detokenize{#2}}}}
    \end{figure}%
  }{%
    \IfFileExists{../../figures/#2}{%
      \begin{figure}[H]
        \centering
        \shotimg{../../figures/#2}%
        \caption{\textbf{#1.} #4\\[0.35em]
          {\sffamily\scriptsize\color{ArduinoComment}#3
            \enspace\textbullet\enspace\texttt{\detokenize{#2}}}}
      \end{figure}%
    }{%
      \noindent
      \begin{tikzpicture}
        \node[
          draw=TblHead, fill=TblSw, line width=1.15pt, rounded corners=6pt,
          inner sep=10pt, text width=\dimexpr\linewidth-22pt\relax, align=left
        ] {
          {\sffamily\bfseries\color{TblHead} Captura #1
            \hfill{\normalfont\ttfamily\scriptsize \detokenize{#2}}}\\[0.35em]
          {\sffamily\scriptsize\color{ArduinoFunc} #3}\\[0.45em]
          {\sffamily\small #4}\\[0.35em]
          {\sffamily\scriptsize\color{ArduinoComment}
            Guardar como \texttt{General\_Documentation/figures/es/\detokenize{#2}} (o \texttt{figures/} compartido).}
        };
      \end{tikzpicture}%
    }%
  }%
  \par\addvspace{0.2em}%
}


\title{Guía de usuario TeleCon4ESP32\\[0.35em]
\large App del teléfono para Panel de control y Vehículo RC}
\author{Para quien usa la aplicación Android\\
\small Autores de firmware: ver el PDF del protocolo de control}
\date{Versión 1.0 \quad 2026}

\begin{document}
\maketitle

\begin{abstract}
Esta guía es para la persona que sostiene el teléfono.
Explica cómo instalar TeleCon4ESP32, cargar un sketch ESP32 de arranque,
conectar y usar \ui{Panel de control} o \ui{Vehículo RC}.
Formatos de bytes, UUID y checksums están en el documento del protocolo.
Pines y banners del Monitor Serie de cada sketch están en la guía de ese sketch.
\end{abstract}

\section{Qué es TeleCon}
\label{sec:what}

El teléfono es siempre el mando.
La placa (ESP32 DevKit, ESP32-CAM o cualquier MCU que hable el contrato)
es la máquina.

Hoy hay dos módulos:

\begin{center}
\renewcommand{\arraystretch}{1.25}
\begin{tabular}{>{\raggedright\arraybackslash}p{3.4cm}
                >{\centering\arraybackslash}p{2.2cm}
                >{\raggedright\arraybackslash}p{6.6cm}}
\rowcolor{TblHead}
\thc{Módulo} & \thc{Acceso} & \thc{Para qué sirve} \\
\rowcolor{TblChan}
\ui{Panel de control} & GRATIS & Sticks, interruptores, knobs, medidores y
gráficas para cualquier proyecto ESP32. \\
\rowcolor{TblSum}
\ui{Vehículo RC} & PRO / monedas & Cabina de conducción y vídeo opcional
de ESP32-CAM. \\
\end{tabular}
\end{center}

\vspace{0.4em}
{\sffamily\scriptsize
\chip{TblChan}{Panel de control --- gratis}
\chip{TblSum}{Vehículo RC --- Pro}
\chip{TblTelem}{El vídeo es un enlace Wi-Fi aparte}}

Los dos módulos comparten la misma sesión de control.
El vídeo (HTTP \texttt{/stream}) es un segundo enlace: únete al Wi-Fi de la
placa para la imagen; Bluetooth o TCP llevan los sticks.

\Needspace{11cm}
\begin{center}
\begin{tikzpicture}[x=1cm, y=0.72cm]
  \fill[rounded corners=7pt, WireBg] (-0.35,-1.2) rectangle (14.1,8.55);
  \draw[rounded corners=7pt, TblHead!28, line width=0.7pt]
    (-0.35,-1.2) rectangle (14.1,8.55);

  \fill[TblHead!11, rounded corners=5pt] (0.1,5.55) rectangle (13.7,6.95);
  \fill[ArduinoPreproc!10, rounded corners=5pt] (0.1,3.55) rectangle (13.7,5.45);
  \fill[ArduinoKeyword!10, rounded corners=5pt] (0.1,1.5) rectangle (13.7,3.45);
  \fill[ArduinoFunc!11, rounded corners=5pt] (0.1,-0.15) rectangle (13.7,1.4);

  \node[telecon actor, fill=TblHead, align=center] at (2.4,7.7)
    {Tel\'efono\\[-1.6pt]{\normalfont\sffamily\scriptsize t\'u}};
  \node[telecon actor, fill=ArduinoPreproc, align=center] at (11.8,7.7)
    {Placa\\[-1.6pt]{\normalfont\sffamily\scriptsize ESP32}};

  \draw[TblHead, line width=1.4pt] (2.4,7.08) -- (2.4,0.12);
  \draw[ArduinoPreproc, line width=1.4pt] (11.8,7.08) -- (11.8,0.12);

  \node[telecon step, fill=TblHead] at (0.55,6.35) {1};
  \draw[telecon arr, TblHead] (2.6,6.5) -- (11.6,6.5);
  \node[telecon pill, fill=white, draw=TblHead!55] at (7.1,6.5)
    {Bluetooth o Wi-Fi de la placa};
  \node[font=\sffamily\scriptsize\bfseries, text=TblHead] at (7.1,5.85)
    {1. Abrir el enlace};

  \node[telecon step, fill=ArduinoPreproc] at (0.55,4.55) {2};
  \draw[telecon arr, ArduinoPreproc] (2.6,5.05) -- (11.6,5.05);
  \node[telecon pill, fill=TblSync] at (7.1,5.05) {El tel\'efono pide empezar};
  \draw[telecon arr, ArduinoKeyword] (11.6,4.2) -- (2.6,4.2);
  \node[telecon pill, fill=TblChan] at (5.35,4.2) {S\'i};
  \node[font=\sffamily\scriptsize\bfseries, text=ArduinoComment] at (7.1,4.2) {o};
  \node[telecon pill, fill=TblId] at (8.85,4.2) {No};
  \node[font=\sffamily\scriptsize\bfseries, text=ArduinoPreproc] at (7.1,3.75)
    {2. La sesi\'on viva empieza solo tras el S\'i};

  \node[telecon step, fill=ArduinoKeyword] at (0.55,2.55) {3};
  \draw[telecon arr, ArduinoKeyword] (2.6,3.05) -- (11.6,3.05);
  \node[telecon pill, fill=TblChan] at (7.1,3.05) {Sticks, interruptores, knobs};
  \draw[telecon arr, ArduinoFunc] (11.6,2.15) -- (2.6,2.15);
  \node[telecon pill, fill=TblTelem] at (7.1,2.15) {Medidores, paneles, gr\'afica};
  \node[font=\sffamily\scriptsize\bfseries, text=ArduinoKeyword] at (7.1,1.7)
    {3. T\'u conduces; la placa puede enviar telemetr\'ia};

  \node[telecon step, fill=ArduinoFunc] at (0.55,0.7) {4};
  \draw[telecon arr, ArduinoFunc] (2.6,0.82) -- (11.6,0.82);
  \node[telecon pill, fill=TblSum] at (7.1,0.82) {Cae el enlace};
  \node[font=\sffamily\scriptsize\bfseries, text=ArduinoFunc] at (7.1,0.18)
    {4. La placa debe parar motores};

  \node[font=\sffamily\scriptsize, text=ArduinoComment] at (7.0,-0.75)
    {Hasta que la placa responda S\'i, mover un stick no es sesi\'on TeleCon viva.};
\end{tikzpicture}
\end{center}

\section{Qué necesitas}
\label{sec:need}

\begin{itemize}[leftmargin=1.4cm]
  \item Un teléfono Android (Android~12 o posterior simplifica el Bluetooth).
  \item Cable USB y un ordenador con Arduino IDE (para cargar el sketch).
  \item Un ESP32 DevKit para la primera sesión.
        Un ESP32-CAM es opcional y solo hace falta para vídeo.
\end{itemize}

Quédate en \ui{Predeterminado} y \ui{Classic + Simple} hasta que ese camino
funcione. \ui{Avanzado} (Binary, BLE, Wi-Fi de DevKit, CAM Binary) es para
después.

\section{Seguridad}
\label{sec:safety}

\begin{itemize}[leftmargin=1.4cm]
  \item Trata los sticks como vivos cuando la app muestre \ui{Conectado}.
        Ruedas, servos y ESC pueden moverse.
  \item Hasta que la placa responda sí, el movimiento en pantalla no es una
        sesión viva aunque Android diga que el Bluetooth está conectado.
  \item Si el teléfono pierde el enlace, los sketches oficiales entran en
        fail-safe (motores parados). Compruébalo en tu propio firmware
        antes de que un vehículo salga del banco.
  \item Un teléfono a la vez. Cierra otras apps que retengan el mismo
        Bluetooth o Wi-Fi.
\end{itemize}

\section{Primera sesión}
\label{sec:first}

Objetivo: en unos diez minutos, mover un stick en \ui{Panel de control} y
ver reaccionar la placa (eco en el Monitor Serie, o un motor si ya hay
cableado).

\paragraph{Antes de empezar.}
Teléfono cargado. ESP32 alimentado por USB. Arduino IDE instalado.
No hace falta Pro ni monedas en este camino.

\subsection*{1 --- Instalar la app}

Instala TeleCon4ESP32 desde Google Play (o tu build de depuración).
Ábrela. Concede Dispositivos cercanos / Bluetooth cuando Android lo pida.
En Android~11 e inferior aún hace falta el permiso de Ubicación para
escanear; la app no necesita que la Ubicación siga activada.

\shot{S01}{s01_home.png}{Vertical, tema claro, no conectado}
{Pantalla Inicio tras el primer arranque. Barra superior (estado
No conectado, PRO dorado opcional, chip de monedas, Ayuda), las dos tarjetas
Panel de control (GRATIS) y Vehículo RC (PRO), e icono de documentos en una
tarjeta. Sin tapar etiquetas. Interfaz en español.}

\subsection*{2 --- Abrir las guías}

En Inicio, toca el icono de documentos de \ui{Panel de control}.
\ui{Codigos y Documentos} muestra la \textbf{Guía rápida} y la
\textbf{Guía de usuario}.
El ZIP del sketch de arranque está en \ui{Firmware correspondiente}
dentro de ajustes (paso siguiente).

\shot{S02}{s02_codes.png}{Vertical, tema claro}
{Códigos y documentos de Panel de control. Solo Guía rápida y Guía de
usuario. Interfaz en español.}

\subsection*{3 --- Obtener el sketch, ajustar y conectar}

Abre \ui{Panel de control} $\rightarrow$ ajustes (engranaje).
\ui{Tipo de usuario} = \ui{Predeterminado}.
\ui{Conexión} = Bluetooth, \ui{Classic + Simple (líneas de texto)}.
En \ui{Firmware correspondiente}, toca el enlace de código de
\texttt{TeleCon\_Classic\_Simple}, comparte o guarda el ZIP, cópialo al
ordenador y abre la carpeta del \texttt{.ino} en Arduino IDE.

\shot{S03}{s03_settings_default.png}{Vertical, tema claro, Tipo de usuario visible}
{Ajustes de Panel de control. Tipo de usuario Predeterminado, Conexión
Bluetooth + Classic + Simple (líneas de texto), y Firmware correspondiente
con el enlace de TeleCon Classic Simple. Interfaz en español.}

Carga \textbf{TeleCon Classic Simple}.
Placa y Puerto en Herramientas deben coincidir con el USB.
Monitor Serie a 115200. Pulsa RESET.
Cuando aparezca el banner de listo, la placa anuncia
\texttt{ESP32-TC-RC-BT-Simple}.

\shot{S02b}{s01_serial.png}{Ordenador, Monitor Serie de Arduino IDE}
{Monitor Serie a 115200 tras RESET. Banner de listo de Classic Simple y
nombre BT \texttt{ESP32-TC-RC-BT-Simple} visibles. Recorta a la ventana Serie.}

Los pines y ganchos de motor están en el \texttt{README.md} de ese sketch,
junto al \texttt{.ino}.

Sal de ajustes y toca el botón Bluetooth en la barra del módulo.

Si Android pide activar Bluetooth, acéptalo.

\shot{S04}{s04_permission.png}{Diálogo del sistema, vertical}
{La hoja real de permiso o ``activar Bluetooth'' que muestra la app
(Dispositivos cercanos en Android 12+, o Ubicación en Android 11 e inferior).
Incluye el nombre de la app en el diálogo.}

En \ui{Conectarse a un dispositivo}, toca actualizar para escanear.
Cuando aparezca \texttt{ESP32-TC-RC-BT-Simple} bajo
\ui{Dispositivos escaneados}, detén el escaneo y toca esa fila.
Si ya estaba vinculado, usa \ui{Dispositivos vinculados}.

\shot{S05}{s05_bluetooth_scan.png}{Vertical, tema claro, escaneo recién detenido}
{Pantalla Bluetooth. Dispositivos vinculados y Dispositivos escaneados
visibles. ESP32-TC-RC-BT-Simple destacado en escaneados.
Control de actualizar / detener visible en la barra superior.}

\paragraph{Resultado esperado.}
La app vuelve al módulo. El estado muestra \ui{Conectado}.
Inicio también muestra Conectado y resalta Panel de control.

\shot{S06}{s06_home_connected.png}{Vertical, tema claro, conectado}
{El mismo Inicio que S01, pero el estado superior dice Conectado y la
tarjeta Panel de control es el módulo reciente resaltado.}

\subsection*{4 --- Abrir Panel de control y mover un stick}

Abre \ui{Panel de control} (horizontal).
Arrastra el stick izquierdo.
En el banco, el firmware Simple oficial refleja el movimiento en medidores
y gráfica (modo simulación, sin cable extra).
En el Monitor Serie deberías ver tráfico de control.

\shot{S07}{s07_control_panel.png}{Horizontal, tema claro o plástico de la app}
{HUD completo del Panel de control: ambos sticks, S1--S6, ambos knobs,
botones laterales, medidores superiores, paneles de siete segmentos, LEDs
y gráfica central. Conectado (sin banner ``No conectado''). Opcional: stick
izquierdo fuera de centro para que se lea el layout.}

Si no se mueve nada, quédate en esta pantalla, confirma Conectado y que
los ajustes siguen en Classic + Simple --- ver Sección~\ref{sec:fix}.

\section{Inicio}
\label{sec:home}

Inicio es el lanzador, no el mando en vivo.

\begin{center}
\renewcommand{\arraystretch}{1.22}
\begin{tabular}{>{\raggedright\arraybackslash}p{3.6cm}
                >{\raggedright\arraybackslash}p{8.8cm}}
\rowcolor{TblHead}
\thc{Control} & \thc{Qué hace} \\
\ui{Panel de control} / \ui{Vehículo RC} &
  Abre ese módulo. Expande una tarjeta para detalles, documentos o
  desbloqueo con monedas. \\
\rowcolor{TblTelem}
Icono de documentos &
  Guía de usuario y Guía rápida de ese módulo
  (los ZIP de sketches están en Ajustes $\rightarrow$ Firmware correspondiente). \\
\ui{PRO} dorado &
  Suscripción / mejora. \\
\rowcolor{TblSw}
Chip de monedas &
  Saldo y precios (si la economía de monedas está activa). \\
Ayuda (interrogación) &
  FAQ corto de Inicio. \\
\rowcolor{TblEnd}
Título \ui{TeleCon4ESP32} &
  Acerca de, privacidad y documentación en línea. \\
\end{tabular}
\end{center}

\ui{Explorar catálogo} solo aparece cuando haya más módulos.
Hasta entonces, toca una tarjeta en Inicio.

\shot{S08}{s08_home_help.png}{Vertical, diálogo de Ayuda abierto}
{Inicio con el diálogo Ayuda rápida. Título Ayuda rápida, al menos dos
temas abiertos (primera configuración y el ESP32 no aparece). Interfaz
en español.}

\section{Conectar}
\label{sec:connect}

Conecta \emph{desde el módulo} que vas a usar (Bluetooth / Conectar de la
barra), no solo desde un emparejamiento viejo del sistema.
Los ajustes del módulo dicen a la app qué enlace abrir.

\subsection{Bluetooth}

Lo usan los sketches de DevKit: Classic Simple (Predeterminado),
Classic Binary, BLE Binary (Avanzado).

\begin{enumerate}[leftmargin=1.5cm]
  \item Concede Dispositivos cercanos (o Ubicación en Android antiguo).
  \item Escanea y elige la placa. El nombre anunciado es solo una etiqueta;
        la app conecta por dirección.
  \item Nombre oficial Predeterminado: \texttt{ESP32-TC-RC-BT-Simple}.
\end{enumerate}

Si un firmware anterior sigue vinculado, Android puede abrir el perfil
equivocado. Olvida el dispositivo en los ajustes Bluetooth del sistema y
vuelve a escanear.

\subsection{Wi-Fi de la placa (SoftAP)}

Lo usan el vídeo CAM y algunos builds Avanzado de DevKit por Wi-Fi.
El teléfono deja el Wi-Fi de casa. Sin internet en esa red es normal.

Red CAM Predeterminada: SSID \texttt{TeleCon-RC-CAM-Starter},
contraseña \texttt{telecon1234}.
Luego toca Conectar en el módulo para que el control use TCP
\texttt{192.168.4.1:3333} (Rol~B) o se quede en Bluetooth (Rol~A).

\shot{S09}{s09_wifi_join.png}{Vertical, ajustes Wi-Fi de Android}
{Selector Wi-Fi del sistema con TeleCon-RC-CAM-Starter (o TeleCon-RC-CAM)
visible, hoja de contraseña o estado ``conectado, sin internet''.
Recorta a la UI de Wi-Fi; TeleCon puede estar en segundo plano.}

No pongas SoftAP y Bluetooth en un solo ESP32-CAM.
CAM de vídeo + DevKit de control es el montaje de dos placas (Rol~A).

\section{Panel de control}
\label{sec:panel}

Mando horizontal para cualquier proyecto ESP32.
Mientras esta pantalla está abierta, sticks y knobs se envían sin parar.

\subsection{Disposición}

Los números coinciden con las llamadas que añadirás sobre la captura~S07:

\begin{enumerate}[leftmargin=1.5cm]
  \item Medidores superiores --- valor analógico y nivel tipo batería.
  \item Paneles de siete segmentos --- números grandes (por ejemplo RPM).
        Títulos y color salen de ajustes; los valores, de la placa.
  \item LEDs de estado --- ocho lámparas según un bitfield de telemetría.
  \item Joysticks --- izquierdo y derecho. Muelle o retención es un ajuste.
  \item Interruptores S1--S6 --- tres por lado.
  \item Knobs --- valores auxiliares tipo 0--100\,\%.
  \item Botones de acción --- cuatro pulsadores laterales.
  \item Centro --- gráfica en Predeterminado; Cámara, Radar y gráfica de
        stick son extras de desbloqueo.
\end{enumerate}

\subsection{Ajustes del panel}

Desde el módulo: ajustes $\rightarrow$ interruptores iniciales,
knobs iniciales, etiquetas de panel, nombres de canales de la gráfica.

En el HUD en vivo del Panel de control, \textbf{toca dos veces el stick}
que quieras configurar (izquierdo o derecho). Se abre la hoja de joystick
de ese stick (Muelle / Retención, orientación del eje, forma).

\shot{S10}{s10_panel_settings.png}{Horizontal, Ajustes de joystick visibles}
{HUD del Panel de control tras un doble toque en un stick: hoja de ajustes
de joystick abierta (Muelle / Retención, orientación del eje, forma).
Interfaz en español.}

La pantalla central (Gráficas / Stick / Cámara / Radar) es un selector
circular en el HUD. Cámara y Radar necesitan monedas o Pro.

\shot{S11}{s11_center_picker.png}{Horizontal, selector de pantalla central abierto}
{Panel de control con la hoja Pantalla central: Gráficas, Stick, Cámara,
Radar. Muestra cuáles están bloqueados si la cuenta es gratuita.}

\section{Vehículo RC}
\label{sec:vehicle}

Cabina Pro.
Se desbloquea con suscripción, monedas o el regalo explorador (anuncio
corto).

Dos montajes de hardware:

\begin{center}
\renewcommand{\arraystretch}{1.25}
\begin{tabular}{>{\raggedright\arraybackslash}p{2.4cm}
                >{\raggedright\arraybackslash}p{4.4cm}
                >{\raggedright\arraybackslash}p{5.4cm}}
\rowcolor{TblHead}
\thc{Rol} & \thc{Placas} & \thc{Qué hace el teléfono} \\
\rowcolor{TblChan}
B Predet. & Un ESP32-CAM &
  Únete al Wi-Fi CAM. Vídeo HTTP. Control por TCP :3333. \\
\rowcolor{TblKnob}
A overlay & CAM + DevKit &
  Wi-Fi CAM solo para vídeo. Sticks por Bluetooth del DevKit. \\
\end{tabular}
\end{center}

\vspace{0.35em}
{\sffamily\scriptsize
\chip{TblChan}{Una CAM --- Predeterminado}
\chip{TblKnob}{Vídeo CAM + Bluetooth DevKit}
\chip{TblId}{No una CAM para SoftAP y Bluetooth a la vez}}

El trim de dirección está en Tune del HUD (tickers $\rightarrow$ Check).
La placa guarda el sesgo.

\shot{S12}{s12_rc_vehicle.jpg}{HUD horizontal}
{Pantalla en vivo de Vehículo RC: dual sticks, telemetría HUD y panel de
cámara (MJPEG en vivo o la tarjeta de ayuda si el stream está inactivo).
Interfaz en español. Preferible Rol B Predeterminado con vídeo en línea.}

\shot{S13}{s13_vehicle_settings.png}{Vertical, sección Conexión / CAM}
{Ajustes de Vehículo RC: Función CAM Un ESP32-CAM (simple), Tipo de usuario
Predeterminado, Rol B SoftAP + Simple, SSID TeleCon-RC-CAM-Starter.
Interfaz en español.}

\section{Predeterminado frente a Avanzado}
\label{sec:modes}

\ui{Tipo de usuario} en los ajustes del módulo oculta a propósito las
opciones difíciles.

\begin{center}
\renewcommand{\arraystretch}{1.22}
\begin{tabular}{>{\raggedright\arraybackslash}p{3.2cm}
                >{\raggedright\arraybackslash}p{9.0cm}}
\rowcolor{TblHead}
\thc{Tipo} & \thc{Enlace típico} \\
\rowcolor{TblChan}
Predeterminado &
  DevKit Classic Simple, o una CAM SoftAP Simple, o Rol~A
  (CAM solo-vídeo + DevKit Classic Simple). Gratis. \\
\rowcolor{TblSync}
Avanzado &
  Classic Binary, BLE Binary, DevKit Wi-Fi Binary, CAM SoftAP Binary.
  Monedas o Pro. \\
\end{tabular}
\end{center}

El firmware de la placa debe coincidir con la fila elegida.
Tras un cambio, desconecta y vuelve a conectar.
\ui{Restaurar valores originales} deja ese módulo en el enlace
Predeterminado de arranque.

\section{Códigos y Acerca de}
\label{sec:docs}

\begin{itemize}[leftmargin=1.4cm]
  \item \ui{Codigos y Documentos} --- esta Guía de usuario y la Guía rápida
        del módulo.
  \item \ui{Ajustes} del módulo $\rightarrow$ \ui{Firmware correspondiente}
        --- comparte o guarda el ZIP del sketch que coincide con la placa y
        la conexión elegidas.
  \item \ui{Acerca de} --- versión, privacidad y el sitio de documentación.
\end{itemize}

\section{Monedas y Pro}
\label{sec:coins}

\ui{Panel de control} sigue gratis en Simple Predeterminado.
\ui{Vehículo RC}, las opciones Avanzado y las vistas extra del centro
(Cámara, Radar, gráfica de stick, CSV de sesión) necesitan monedas o
suscripción Pro.

\paragraph{Dónde mirar.}
El chip de monedas en la barra de Inicio muestra el saldo.
Tócalo para abrir \ui{Precios en monedas} (monedero y tarifa).
En una tarjeta Pro bloqueada usa las acciones de desbloqueo / monedas, o
\ui{Hazte Pro} al pie de Inicio para una suscripción que evita monedas y
anuncios.

\paragraph{Ganar monedas.}
Mira un anuncio recompensado desde la hoja de precios (botón de reproducción,
normalmente \texttt{+3} monedas por anuncio). A veces el anuncio no está
disponible; inténtalo más tarde.

\paragraph{Gastar monedas.}
Las monedas desbloquean \textbf{un módulo Pro a la vez} durante el plazo
que elijas. Opciones típicas (el coste exacto está en la hoja de la app):

\begin{center}
\renewcommand{\arraystretch}{1.15}
\begin{tabular}{>{\raggedright\arraybackslash}p{3.2cm}
                >{\raggedright\arraybackslash}p{8.2cm}}
\rowcolor{TblHead}
\thc{Opción} & \thc{Significado} \\
Un uso & Desbloqueo de una sola sesión de ese módulo. \\
\rowcolor{TblTelem}
24 horas & Acceso temporizado de un día. \\
3 días / 1 semana & Acceso más largo; suele mejorar el coste por día. \\
\end{tabular}
\end{center}

Abre el diálogo de desbloqueo en el módulo deseado, elige la duración y
confirma si el saldo alcanza. Los paquetes largos suelen bajar el
``por día'' de la tabla de precios.

\shot{S15}{s15_coin_pricing.png}{Vertical, hoja Precios en monedas}
{Diálogo de monedero / Precios en monedas: tasa de ganancia (+3 por anuncio
recompensado), texto de gasto y tabla Opción / Coste / Usos / Por día.
Interfaz en español.}

\paragraph{Regalo explorador.}
El brillo dorado en una tarjeta bloqueada es un regalo explorador ocasional:
mira un anuncio corto y desbloqueas ese módulo Pro unas horas (unas cuatro).
El brillo vuelve al cabo de unos días.

\paragraph{Suscripción Pro.}
\ui{Hazte Pro} desbloquea los módulos Pro sin gastar monedas ni el ciclo de
anuncios. Detalle y precio están en la hoja de suscripción de Play.

\shot{S14}{s14_coin_unlock.png}{Vertical, diálogo de desbloqueo}
{Diálogo de desbloqueo con monedas para un módulo concreto (por ejemplo
Vehículo RC o ajustes Avanzado): plazos y confirmación. Interfaz en español.}

\section{Cuando falla}
\label{sec:fix}

Empieza por el síntoma, no por el protocolo.

\begin{center}
\renewcommand{\arraystretch}{1.2}
\begin{tabular}{>{\raggedright\arraybackslash}p{3.7cm}
                >{\raggedright\arraybackslash}p{8.5cm}}
\rowcolor{TblHead}
\thc{Ves} & \thc{Prueba esto} \\
La placa no sale al escanear &
  Alimentación + RESET. Banner de listo a 115200.
  Dispositivos cercanos concedidos. Acércate. Detén el escaneo al verla. \\
\rowcolor{TblTelem}
Falla al conectar &
  Abre Bluetooth desde el módulo, no solo desde el sistema.
  Olvida un vínculo viejo. Classic frente a BLE igual que el sketch. \\
Inicio sigue No conectado &
  Conecta dentro del módulo. Un solo teléfono. Luego Inicio se actualiza. \\
\rowcolor{TblSw}
Sticks / medidores muertos &
  Quédate en la pantalla en vivo. Los ajustes deben coincidir con el
  sketch (Classic + Simple Predeterminado en la primera sesión).
  Reflashea si el Serial nunca muestra listo. \\
No se pudo conectar /
  tiempo de espera &
  El teléfono abrió el tubo pero la placa no respondió a tiempo.
  El mismo modo a ambos lados; actualiza el sketch si ignora el arranque. \\
\rowcolor{TblId}
El protocolo no coincide &
  Ajustes en Simple y la placa habla Binary (o al revés).
  Cambia ajustes o carga el ZIP correcto. \\
Sin cámara &
  Únete primero al SSID CAM. Contraseña \texttt{telecon1234}.
  Rol~A: vídeo en Wi-Fi CAM, sticks en Bluetooth del DevKit --- no ambas
  radios en una CAM. \\
\rowcolor{TblSum}
Los motores siguen
  tras la caída &
  Los sketches oficiales paran. Si escribiste firmware, añade fail-safe
  antes de probar un vehículo. \\
\end{tabular}
\end{center}

Xiaomi / POCO en SoftAP: apaga datos móviles si muere el TCP, y deja
TeleCon en primer plano. VPN desactivada.

\section{Adónde ir después}
\label{sec:next}

\begin{itemize}[leftmargin=1.4cm]
  \item \textbf{README} del sketch junto al \texttt{.ino} --- pines, SSID,
        nombre BT, \texttt{YOUR CODE HERE}.
  \item PDF del \textbf{protocolo de control} --- si escribes tu propio firmware.
  \item Ayuda en la app --- recordatorios cortos en el teléfono.
\end{itemize}

\end{document}
